% Unit 2 map -- one page.
\documentclass{apchem-worksheet}

\apcunit{2}
\apcunittitle{Compound Structure and Properties}
\apcdoctitle{Unit Map}
\apcsubtitle{CED Unit 2 \textbullet\ 7--9\% of the AP Exam \textbullet\ 6 blocks + exam}

\begin{document}
\apctitleblock

\begin{apcnote}[How this unit is graded]
Unit exam Wed Sep 9: 20 multiple choice (\textbf{no calculator}) + 1 long
free-response (10 points), 52 minutes. This unit is where AP stops asking
\emph{what} and starts asking \emph{draw it, then defend it} --- the FRQ
rewards justified structures, not memorized shapes.
\end{apcnote}

\section*{\sffamily Block calendar}

\renewcommand{\arraystretch}{1.35}
\noindent\begin{tabular}{@{}p{2.2cm}p{5.2cm}p{1.4cm}p{3.6cm}p{2.0cm}@{}}
\toprule
\textbf{Date} & \textbf{Topics} & \textbf{CED} & \textbf{Read (PDF pp.)} & \textbf{Practice}\\
\midrule
Wed Aug 26 & Bond types; electronegativity; potential-energy curves & 2.1, 2.2 & \S8.1--8.3 \textbullet\ 391--400; \S8.8 \textbullet\ 412--415 & WS 1\\
Fri Aug 28 & Ionic solids \& lattice energy; metals and alloys & 2.3, 2.4 & \S8.5 \textbullet\ 404--408; \S10.4, 10.7 \textbullet\ 505--511, 520--523 & WS 1\\
Mon Aug 31 & Lewis diagrams & 2.5 & \S8.9--8.11 \textbullet\ 415--423 & WS 2\\
Wed Sep 2 & Formal charge; resonance & 2.6 & \S8.12 \textbullet\ 423--427 & WS 2\\
Fri Sep 4 & VSEPR: electron domains and molecular geometry & 2.7 & \S8.13 \textbullet\ 427--440 & WS 3\\
Mon Sep 7$^{*}$ & Hybridization; molecular polarity; synthesis & 2.7 & \S9.1 \textbullet\ 455--466 & WS 4 (FRQ)\\
\textbf{Wed Sep 9} & \textbf{UNIT 2 EXAM} & all & review sheet & ---\\
\bottomrule
\end{tabular}

{\sffamily\footnotesize $^{*}$Mon Sep 7, 2026 is Labor Day --- confirm your
co-op meets; if not, everything slides one block and the exam moves to
Fri Sep 11.}

\vspace{0.5em}
\section*{\sffamily By the end of this unit you can\ldots}

\begin{itemize}[leftmargin=*,itemsep=0.15em]
  \item Classify bonds as ionic, polar covalent, nonpolar covalent, or
        metallic from electronegativity and position --- and say why the
        categories blur. \ced{2.1}
  \item Read a potential-energy vs.\ distance curve: bond length at the
        minimum, bond energy as the well depth. \ced{2.2}
  \item Rank lattice energies with Coulomb's law (charges first, then
        radii). \ced{2.3}
  \item Explain metallic bonding and predict how interstitial and
        substitutional alloys change properties. \ced{2.4}
  \item Draw Lewis diagrams fast and correctly, including the octet
        exceptions worth knowing. \ced{2.5}
  \item Use formal charge to choose between candidate structures; draw
        resonance sets and state what they mean for bond length. \ced{2.6}
  \item Predict electron-domain and molecular geometry, bond angles,
        hybridization (sp, sp$^2$, sp$^3$), and molecular polarity. \ced{2.7}
\end{itemize}

\begin{apctrap}[Where students lose points in this unit]
Drawing one resonance form and calling it ``the'' structure; skipping the
formal-charge check that picks the best skeleton; claiming a molecule is
polar because it ``has polar bonds'' (symmetric \ce{CO2} says otherwise);
quoting 109.5$^\circ$ for bent molecules without the lone-pair compression
argument; and explaining lattice energy by size when the charges differ ---
charges dominate.
\end{apctrap}

\end{document}
